\section{\label{III_1_A}Une segmentation manuelle} 

La préparation des pratiques utilisateurs à ces nouvelles fonctionnalités passe notamment par une préparation méthodologique de l'interface, dédiée aux collections audiovisuelles. Par exemple, la mise en place d'un outil de visionnage et d'annotation vidéo pourrait permettre de préparer les utilisateurs à de telles pratiques. Cet outil pourrait répondre à un double enjeu, celui d'encourager les pratiques d'annotation et de segmentation sur les fichers numériques, ainsi que celui d'éventuellement collecter les données produites par un tel outil afin d'entraîner les modèles à implémenter. (collection as data)
La mise en place de cet outil permettrait également de faciliter le contrôle qualité, en comparant les données manuelles et automatiques. 

L'outil Mediascope, développé par l'\gls{ina} à destination des chercheurs à l'INAthèque, est un outil d'aide à l'analyse des contenus audiovisuels. Il permet de structurer le développement temporel du document à l'aide d'annotations associées à des "time-codes" et à des images fixes \footcite{zotero-item-746}. L'utilisateur peut ainsi visionner un document audiovisuel, annoter des marqueurs temporels sur des images fixes, ainsi qu'y saisir du texte. L'outil permet également l'export des images, du texte, et de statistiques sur le contenu \footcite{zotero-item-746}. Mediascope unifie ainsi la lecture vidéo et l'indexation d'annotations au sein d'une même interface.

Cette proposition s'inscrit dans un document global d'améliorations : cf annexe :

Au sein de l'application Skinsoft, la gestion des médias se limite à la visualisation simple d'un fichier mutimédia lié à une notice. Afin d'améliorer cette visualisation des médias, nous avons proposé au cours de notre stage, sur l'insipration du Médiascope, une visionneuse permettant le séquençage manuel et l'annotation de données. 
Cette amélioration s'appuierait d'une part sur deux possibilités de séquençage manuel. L’utilisateur peut saisir des repères temporels sur un fichier vidéo sans pour autant le modifier. Ces repères sont reversés dans des champs prévus pour accueillir les “time codes” dans la notice de l’objet numérique. Au sein du fichier vidéo, les marqueurs sont matérialisés et visibles sur la barre de défilement de la vidéo, permettant une navigation interactive. De plus, la visionneuse permettrait à l’utilisateur de segmenter le fichier vidéo en plusieurs séquences distinctes, à partir de repères temporels placés (développement en annexe). Chaque séquence identifiée est reversée dans un fichier d'item numérique propre, lié à une notice descriptive du niveau Fragment. 
D'autre part, cette amélioration permettrait deux niveaux d'annotations sur le média. Une annotation globale au niveau de la séquence qui prendrait par exemple la forme d'un résumé, de mots-clés généraux thématiques, inscrits pau sein de marqueurs temporels. Une annotation précise de tags apposés sur des éléments de l'image tels que des objets, personnes, éléments de décor, similaire à la détection d'objet prévue par \gls{yolo}. Ces éléments seraient indexés dans un champ dédié avec un ancrage temporel et spatial. (développement en annexe)
À la manière du médiascope, cette amélioration du visionnement des médias permettrait de combiner visualisation, segmentation et annotation au sein d'un même outil. 

La mise en place d'un tel outil permettrait déjà d'améliorer la recherche en la faisant porter sur des éléments précis portés par le contenu des images. De plus, le séquençage et l'annotation manuelle sur les documents audiovisuels prépare es utilisateurs à la structure des données engendrées en masse par les algorithmes à implémenter. Enfin, la production manuelle de ces données (time codes et mots-clés) faciliterait l'entraînement et le réentraînement des modèles, ainsi que le contrôle qualité effectué par l'utilisateur. Les données produites en amont permettraient en effet de constituer un jeu de données d'entraînement pour les modèles implémentés. Ces données serviraient également de référence afin d'inspecter la qualité de celles produites automatiquement par les modèles.